\section{Limiti attuali e aree di miglioramento: da proposta a sistema attuato}\label{sec:ps}

Questo capitolo dice, per ogni limite trovato, che cosa manca per passare dalla proposta a un sistema attuato: la miglioria (che cosa e come) e la prova che la decide, con la soglia fissata prima della prova e il risultato che la smentirebbe. Nessuna delle migliorie è ancora provata su un sensore fisico o su hardware dedicato.

\paragraph{Che cosa si può dire dei risultati.} Il prototipo è stato costruito e provato in un giorno. Le prove svolte non coprono tutti i casi, e non potrebbero: per lo più una geometria per scenario, sensori simulati dalle schede, etichette a mano, nessun sensore fisico. Dentro questi limiti i risultati sono incoraggianti, e ciascuno ha una base matematica e un intervallo: nel volo veloce il principio di sicurezza (nessun «libero» sull'oggetto o sui frammenti, con la mappa dichiarata) regge in tutte le celle simulate, $0/200$ $[0;\,0{,}019]$ nel controllo e $0/60$ $[0;\,0{,}060]$ nelle altre; l'insieme di presenza ha contenuto il vero dopo l'urto e l'avaria; a bordo di un drone in moto la regione del drone seguito non si è persa in $60/60$ passaggi $[0{,}940;\,1]$ per geometria con il radar a $360^\circ$, anche durante una collisione vicina; dopo la taratura, su traiettorie di volo pubbliche con un sensore simulato, la regione al $90\,\%$ ha contenuto l'aereo da 85 a 89 volte su 100 fra $1$ e $30\unita{s}$, con i casi concentrati su una rotta e la virata come limite ($0{,}73$ $[0{,}67;\,0{,}78]$). Le cifre sotto soglia sono riportate come limiti, e la tabella che segue dice, per ciascuno, quale miglioria serve e come la si decide.

\subsection{Tabella dei limiti e delle migliorie}\label{sec:ps-tabella}

Le latenze sono stime da ipotesi; nessuna è misurata. «Applicata in simulazione» vuol dire che la miglioria è stata scritta e provata nel simulatore, non su un sensore.

{\footnotesize\setlength{\tabcolsep}{4pt}
\begin{longtable}{@{}>{\RaggedRight\arraybackslash}p{3.3cm}>{\RaggedRight\arraybackslash}p{5.6cm}>{\RaggedRight\arraybackslash}p{5.8cm}@{}}
\caption{Limiti, migliorie e prove che le decidono.}\label{tab:ps-migliorie}\\
\toprule
Limite & Miglioria: che cosa e come & Stato, esito e prova che la decide\\
\midrule
\endfirsthead
\toprule
Limite & Miglioria: che cosa e come & Stato, esito e prova che la decide\\
\midrule
\endhead
\bottomrule
\endfoot
collisione veloce (F3) non dichiarata come evento & regola di contatto: tracce risolte e mature, convergenza $\ge10\unita{m/s}$, fusione o sparizione, tracce nuove o Doppler fuori norma & \emph{applicata in simulazione, non adottata}: B $57/60$ $[0{,}863;\,0{,}983]$, ma falsi eventi in F5 $2/60$ ed eventi prima dell'urto; A $34/60$. Prossimi passi: Doppler sulla cella unita, riagganci esclusi, probabilità di collisione raggiunta, radar vicino su una griglia di geometrie fissata prima. Soglia: F3 in B $\ge54/60$ con F5, F1 e F6 a 0, su semi nuovi e su almeno 10 geometrie\\
insieme di presenza a $90\unita{m/s}$ nel nucleo & inviluppo a $210\unita{m/s}$ anche nell'insieme di presenza & \emph{applicata in simulazione}: F4 in A da $13/60$ a $60/60$; F6 in A da $126/200$ a $186/200$\\
simulatore dei frammenti con l'integratore di Eulero, instabile & integratore semi-implicito & \emph{applicata in simulazione}: collisione fra multirotori, volume da $48/60$ a $60/60$ $[0{,}940;\,1]$; evento $48/60$ e 3 eventi precoci, ora limite\\
confidenza ottimista a 1--$2\unita{s}$ & $\sigma_{\mathrm{rif}}^2$ ($20\unita{m}$ per gli aerei), regime di accelerazione, pesi fissati, ricalibrazione isotonica, minimo progressivo & \emph{applicata su dati pubblici}: da $0{,}71$ a $0{,}89$ $[0{,}87;\,0{,}91]$, scarto medio $0{,}016$. Peggiora la virata (da $0{,}91$ a $0{,}73$), l'accelerazione resta a $0{,}76$. Prossimo: ricalibrazione condizionata a una manovra rilevata in modo causale (soglia: virata $\ge0{,}85$; smentita sotto $0{,}80$), cinque divisioni, divisione per rotta\\
confidenza del nucleo mal tarata con il radar a $1\unita{Hz}$ & minimo progressivo e covarianza gonfiata sui registri di F1 ed F6 & non fatto; soglia: con $C\ge0{,}9$ copertura osservata $\ge0{,}80$\\
avaria non contenuta: $0/10$ $[0;\,0{,}278]$ & profilo dei frammenti con $\Delta v\ge300\unita{m/s}$, impronta $\ge10\unita{km}$, Doppler sui frammenti & non fatto; soglia: $60/60$ passaggi con tutti i frammenti dentro $V$\\
innesco della fiammata ottimista (anche in F7) & saturazione rilevata dall'istogramma della camera & non fatto; soglia: ritardo $\le1$ ciclo in almeno $0{,}9$ dei casi, su registrazioni vere\\
identità dopo la cecità (scambio $0{,}47$--$0{,}63$) & Doppler, sezione radar e firma termica per riassociare & non fatto; soglia: scambio $\le0{,}2$\\
continuità con il radar rotante a $1\unita{Hz}$ & radar a scansione elettronica a $10\unita{Hz}$ (tracce a $10\unita{Hz}$ da stampa, campo di $120^\circ\times80^\circ$ dal produttore; portate per bersaglio non verificate) & simulato in B: $60/60$; manca la misura sul sensore\\
F7: radar solo frontale ($\pm60^\circ$) & radar a $360^\circ$ (quattro pannelli) o camera su brandeggio puntata dalla traccia & \emph{applicata in simulazione}: in B identità in virata $60/60$, evento di lato e dietro $60/60$ (contro $43/60$, $1/60$, $0/60$)\\
F7: radar non accecato dall'esplosione & ricevitore saturo per $0{,}2$--$2\unita{s}$ e occultamento (variante F7r, modello pessimista) & \emph{applicata in simulazione}: «libero» $0/60$ e continuità $59$--$60/60$ reggono; l'avviso prima dei frammenti crolla (cella totale $0/59$); la riacquisizione con il radar frontale scende a $41$--$43/60$. Correzione c1 (avviso immediato «evento vicino non localizzato» alla saturazione congiunta di radar e camere) \emph{applicata in simulazione}: avviso $59/59$ $[0{,}939;\,1]$ nella cella totale; falsi $11/200$ $[0{,}031;\,0{,}096]$ con saturazione del radar d'altra causa; riacquisizione davanti $49/60$, sotto la soglia di $51/60$; non adottata. Prossimo: localizzare l'avviso (direzione del lampo, Doppler); riacquisizione presunta dentro il volume con coerenza Doppler; modello del controllo del guadagno dalla scheda. Soglie: falsi $\le0{,}05$ con avviso $\ge0{,}9$; riacquisizione davanti, radar frontale, $\ge51/60$\\
F7: volume prudente, avviso non discriminante ($60/60$ senza frammenti vicini) & velocità per frammento dal Doppler; volume per grappoli tracciati & non fatto; soglia: avviso falso $\le10\,\%$ con contenimento al $100\,\%$\\
F7: margine dell'avviso sotto $T_r$ con l'urto a 10--$30\unita{m}$ (7 su 20) & previsione dell'urto dalle tracce prima del contatto (avviso di prossimità, come in F3) & non fatto; soglia: margine $\ge0{,}25\unita{s}$ in $20/20$ a 10--$30\unita{m}$\\
F7: velocità vera del portatore nel tracciatore & velocità dalla navigazione inerziale; covarianza della navigazione correlata nello stato & non fatto; stesse soglie della tabella~\ref{tab:vo-scorta}\\
F7: controllo senza disturbo a grappoli & controllo con stormi, pioggia e contromisure simulati & non fatto; soglia: falsi eventi $0/200$\\
F7: tracciatore nuovo, non il nucleo & integrare F7 nel nucleo con il sensore in moto & non fatto\\
una geometria per scenario & almeno 10 geometrie per scenario & in parte: F3 su 5 varianti (16--19 su 20 in B); F7 su 3 varianti\\
tempestività verso l'umano & avviso con la direzione: $5{,}4\unita{s}$ invece di $12{,}5\unita{s}$ (AC 90-48E, Tabella 1 \cite{faa-ac-90-48e}); $2005\unita{m}$ invece di $3312\unita{m}$ a $250\unita{m/s}$ & calcolato; manca la prova con operatori\\
latenza del calcolo & FPGA o ASIC per segnale e CFAR (CFAR su FPGA in $0{,}6\unita{\mu s}$ \cite{sim2023}; distanza-Doppler 85 volte più rapida che su ARM \cite{choi2021}); calcolo dedicato totale $0{,}08$--$1{,}2\unita{ms}$ (ipotesi). Analogia con i chip per modelli linguistici (circa 17\,000 contro circa 230 \emph{token} al secondo, dato aziendale da prove interne \cite{heise2026-taalas}), valida solo per i carichi limitati dalla memoria. Passo di Kalman su FPGA: più lento in latenza, circa $4{,}5$ volte in portata \cite{summers2018} & misura su FPGA reale; soglia: $\le1{,}2\unita{ms}$ per aggiornamento\\
primo avviso verso una macchina & la leva è la rivisita: a $10\unita{Hz}$ $0{,}62\unita{s}$, a $40\unita{Hz}$ $0{,}17\unita{s}$ (calcolato da ipotesi) & misura dalla cattura all'avviso\\
mappa, classificatore, sensore fisico & LiDAR per la mappa, classificatore per classe, prove sul campo & tutte le soglie precedenti\\
\end{longtable}}

\subsection{Piano a tappe}\label{sec:ps-piano}

Ogni tappa riusa le soglie della tabella~\ref{tab:ps-migliorie} e aggiunge una sola fonte di realtà alla volta, in modo che un fallimento si possa attribuire.
\begin{enumerate}[label=\arabic*.]
\item \textbf{Simulazione estesa.} Geometrie variate (almeno 10 per scenario), disturbo a grappoli e uccelli simulati, i passi successivi della regola di contatto e della taratura della confidenza, i frammenti di un'avaria non contenuta, il portatore in manovra e il vento in F7, il tracciatore di F7 nel nucleo. Uscita: le stesse classi della tabella~\ref{tab:vo-risultati}, con le soglie fissate prima. Non richiede hardware.
\item \textbf{Banco con registrazioni.} Il nucleo, compilato, gira su registrazioni vere di radar e camere (archivi pubblici o di terzi, con licenza), senza modificarle. Si misurano mediana, 99° percentile e massimo della latenza di calcolo, l'innesco della fiammata dal segnale della camera, la confidenza contro una verità di bordo più pulita dell'ADS-B, i falsi eventi con il disturbo reale.
\item \textbf{Sensore su banco.} Un radar a scansione elettronica e una camera, fermi, su oggetti noti (droni in volo controllato in un'area autorizzata): rivisita, portata sull'oggetto piccolo, sezione radar dei frammenti simulati con oggetti di prova, latenza dalla cattura all'avviso. Qui si misura per la prima volta la colonna «dedicato» della tabella~\ref{tab:vo-latenza}, su FPGA.
\item \textbf{Campo.} Prove in un sito autorizzato con traffico vero, con operatori: tempestività verso l'umano, falsi avvisi all'ora, mappa verificata con il LiDAR, classificatore per classe. Solo dopo questa tappa una cifra del sistema si potrebbe dire misurata su un sensore.
\end{enumerate}

\subsection{Che cosa resta fuori}\label{sec:ps-fuori}

Il sistema resta un sistema di osservazione e di avviso: non decide manovre, non comanda mezzi, non ha funzioni d'ingaggio o di contrasto. Non sostituisce le procedure di separazione del traffico aereo né i sistemi anticollisione di bordo. Le prove descritte in questo documento non danno alcuna base per un uso operativo senza le tappe 2--4 del piano. Nello scenario F7 il sistema mantiene la traccia di un drone e avvisa il portatore: non lo guida e non agisce sul drone seguito.
